\section{Methods}

\subsection{Problem and notation}
Let $\mathcal{D}=\{(g_i,d_i,T_i,t_i)\}_{i=1}^{n}$ be the labeled genomes, with maximal doubling time $d_i$ (hours) at the
optimal growth temperature $T_i$ and GTDB taxonomy $t_i$. From the sequence of $g_i$ alone we compute an input vector
$\mathbf{x}_i$ (Section~\ref{sec:features}) and learn three predictors: of $y_i=\log d_i$, of $T_i$, and of
$y^{\mathrm{ref}}_i=\log d^{\mathrm{ref}}_i$, the doubling time moved to $T_{\mathrm{ref}}=\TRefC$\,\textdegree C with an
Arrhenius factor,
\begin{equation}
\mu^{\mathrm{ref}}=\mu\exp\!\left[-\frac{E_a}{R}\left(\frac{1}{T_{\mathrm{ref}}}-\frac{1}{T}\right)\right],\qquad \mu=\frac{\ln 2}{d},
\end{equation}
with $E_a=\EaKJ$\,kJ\,mol$^{-1}$. We judge models mainly on $y$, because the correction multiplies the doubling time of a
hot-growing organism by a large factor and a model can look accurate on $y^{\mathrm{ref}}$ partly by recognizing that an
organism grows hot. For a new genome $g$, its novelty level $\ell(g)$ is the lowest rank (genus, family, order) at which
it shares a taxon with a labeled genome, or ``new order''. Each prediction is reported with the error and range measured
at its level. Accuracy is summarized by the typical fold error, the Spearman correlation $\rho$, and the area under the
receiver operating curve (AUC) for $d>24$\,h:
\begin{equation}
E=\exp\Big(\operatorname*{median}_i\,|y_i-\hat y_i|\Big),\qquad
\hat y_g\pm=\big[\hat y_g+q^{\ell,\delta}_{0.1},\ \hat y_g+q^{\ell,\delta}_{0.9}\big],
\label{eq:metrics}
\end{equation}
where $q^{\ell,\delta}_{p}$ is the $p$-quantile of the nested held-out residuals $y_i-\hat y_i$ at level $\ell$ in
domain $\delta$ (bacteria or archaea), which gives an 80\% range. Its coverage is checked by cross-fitting: the quantiles
are estimated on four fifths of the genera and applied to the rest. Because a range cannot target organisms selected by
their true doubling time, each genome also gets a probability of $d>24$\,h from a logistic model on $\hat y$ and the
domain, cross-fitted in the same way.

\subsection{Training labels and their audit}
\paragraph{Labels.} The core is the growth-rate table distributed with gRodon (\NGrowthTable species, derived from the
compilation of \citet{madin2020}), used in three ways: species whose genome we placed and processed ourselves; species for
which only the trait table gives genome size, GC content and tRNA count; and species with a genome accession but no traits,
for which every input was computed from that genome. We added \NLiteratureWord strains from lineages the compilation lacks,
each checked in the primary paper and paired with the genome of the measured strain
\citep{konneke2005,stieglmeier2014,palatinszky2015,nowka2015,mohammadi2017,lehtovirtamorley2014}. Names are mapped to
GTDB R\GTDBRelease \citep{parks2022}, and rows are joined into species groups (any two names sharing a GTDB species, and
any two GTDB species sharing a name), so that synonyms never fall on both sides of a split. One genome is listed under two spellings of its species; both rows fall in one species group, and dropping one changes
no headline error or correlation by more than \DupMaxChange.

\paragraph{Audit.} A label is meant to be a maximal rate at the stated temperature. We screened labels at three levels.
(i) Species: a label at least \ScreenFoldWord-fold from both its held-out prediction and the median of its labeled genus
relatives was traced to its source records. (ii) Study (Algorithm~\ref{alg:audit}): because one study can shift a whole
genus and leave no outlier among relatives, every label was traced to its study and each study was compared with evidence
from other studies only. (iii) Records: labels matching no source record but far from the closest one were listed, to
catch unit slips. A flag is not a removal: a label was removed only when a second source or the study itself showed that
it was not a maximal rate (Supplementary Table~\ref{S-tab:exclusions}).

\begin{algorithm}[t]
\caption{Study-level label check}\label{alg:audit}
\footnotesize
\begin{algorithmic}[1]
\Require labels $d_i$ with species; source records $(v_r,\mathrm{study}_r)$ of the compilation
\For{each label $i$}
  \State $s_i\gets$ study of the record with $|v_r-d_i|\le\TraceTolPct\%\cdot d_i$, if any
  \State $e_i\gets$ median of other studies' values for the species, else median label of $\ge 2$ genus (or family)
  relatives labeled by other studies
\EndFor
\For{each study $s$ with $\ge 2$ labels that have evidence}
  \State $r_i\gets\log(d_i/e_i)$ for its labels; $\tilde r\gets\operatorname{median}_i r_i$
  \If{$|\tilde r|\ge\log\StudyFold$ \textbf{and} $\ge 80\%$ of $r_i$ share the sign of $\tilde r$} flag $s$ \EndIf
\EndFor
\State list untraced labels $\ge\StudyFold$-fold from the closest record (unit or conversion slips)
\State \Return flags, each reviewed against a second source before any removal
\end{algorithmic}
\end{algorithm}

\subsection{Genome inputs}\label{sec:features}
All inputs come from the genome sequence, so they are computed in the same way for an isolate and for a
metagenome-assembled genome.

\paragraph{Codon usage.} Genes are predicted with Prodigal through pyrodigal \citep{hyatt2010,larralde2022} on contigs of
at least \MinContig\,bp, trained on the genome when it is at least \MinTrainKb\,kb long and in metagenomic mode otherwise.
Ribosomal protein genes, the usual proxy for highly expressed genes, are found with \NRibosomalFamilies Pfam families
\citep{mistry2021} using HMMER \citep{eddy2011} and the gathering thresholds. Codon usage bias is the median MILC distance
\citep{supek2005} of the ribosomal protein genes from the codon usage of the whole genome, required to rest on at least
\MinRibosomal genes. On the same genes we run the released gRodon \citep{weissman2021} without a temperature and keep its
\NGRodonFeaturesWord statistics (codon usage bias of the highly expressed genes, its consistency, codon pair bias and gRodon's
own doubling time). The share of the amino acids IVYWREL in all proteins, which rises with growth temperature
\citep{zeldovich2007}, signals the temperature.

\paragraph{Lifestyle genes.} The model has \NLifestyleFeatures inputs from \NLifestyleFamilies Pfam families with their
gathering thresholds (Supplementary Table~\ref{S-tab:pfam}) that describe how a genome makes a living: terminal oxidases, complex I,
hydrogenases, methyl-coenzyme M reductase, DsrD, nitrogenase, RuBisCO, ammonia or methane monooxygenase, photosynthetic
reaction centres, catalase, flagellin, and the densities of transporter and regulator genes per thousand genes. Each call
was checked against well-known organisms before use.

\paragraph{Phylogeny and 16S.} A taxon's place on the tree is its vector of distances to \NLandmarks landmark tips of the
GTDB R\GTDBRelease bacterial or archaeal tree, chosen once by farthest-point sampling, so that the vector of a genome does
not depend on which other genomes are analyzed; a genome takes the mean vector of the tips of its GTDB species, or of its
genus, family or order when the species is not on the tree. The 16S rRNA gene is described in the same way, by cosine
distances between \KmerK-mer profiles and \NLandmarks reference sequences. The 16S gene is cut from the genome with barrnap
\citep{barrnap} for metagenome-assembled and labeled genomes alike (\SixteenNGenomeCopy of the \NGenomes labeled genomes);
a labeled genome without a full-length copy uses a 16S sequence of its species from NCBI. The other inputs are genome size,
GC content, tRNA genes (tRNAscan-SE \citep{chan2021}) and a flag for archaea.

\subsection{Model}
Each target is predicted by
\begin{equation}
f(\mathbf{x})=\tfrac12\big[f_{H}(\mathbf{x};\theta_H)+f_{R}(\mathbf{x};\theta_R)\big],
\end{equation}
the average of histogram gradient boosting with absolute (median) loss \citep{friedman2001} and a random forest of
\NForestTrees trees \citep{breiman2001} (scikit-learn, \citealp{pedregosa2011}). The settings $\theta_H$ and $\theta_R$
are chosen separately, each as the best of \NTuneCandidates random draws from a grid (Supplementary
Table~\ref{S-tab:settings}) by typical error in a \NTuneFolds-fold split grouped by GTDB family, averaged over
\NTuneShuffles shuffles. Most new genomes are placed on the tree only to genus or a coarser rank, while a labeled genome is
placed at its species; each labeled genome therefore enters training four times,
$\mathcal{D}^{\mathrm{aug}}=\bigcup_{r}\{(\mathbf{x}^{(r)}_i,y_i)\}$ for $r\in\{\text{species},\text{genus},\text{family},\text{order}\}$,
where $\mathbf{x}^{(r)}_i$ carries the tree vector of its taxon at rank $r$. A genome lacking a 16S gene or a place on the
tree is predicted by the model trained without that block.

\subsection{Evaluation}
\paragraph{Lineage-held-out splits.} A random split puts close relatives of every test genome in training, which most
environmental genomes do not have. We therefore hold out whole species groups, genera, families or orders
(Algorithm~\ref{alg:nested}) among the \NHeadlineRows labeled genomes whose genes could be predicted, and the whole
procedure of the previous section, including the choice of settings, is
repeated inside every split, so that held-out genomes take no part in any choice. The held-out error at level $\ell$ is the
error attached to a new genome whose nearest labeled relative sits at that level. For comparison we also run the model
with the settings chosen on all labeled genomes and each genome entered once, the scheme behind the error table released
with the predictions.

\begin{algorithm}[t]
\caption{Nested lineage-held-out evaluation}\label{alg:nested}
\footnotesize
\begin{algorithmic}[1]
\Require labeled set $\mathcal{D}$; held-out rank $L\in\{$species group, genus, family, order$\}$; seeds $s=1..\NSeedsHeadline$
\For{each seed $s$}
  \State shuffle the taxa at rank $L$ with seed $s$; split them into \NFolds folds (no taxon in two folds)
  \For{each fold $k$ with training rows $\mathcal{T}$ and test rows $\mathcal{V}$}
    \State $\theta_H,\theta_R\gets$ settings search on $\mathcal{T}$ only (as in the deployed model)
    \State fit $f$ on $\mathcal{T}^{\mathrm{aug}}$ (four placements per genome)
    \State $\hat y_i\gets f(\mathbf{x}^{(\text{species})}_i)$ for $i\in\mathcal{V}$
  \EndFor
  \State $E_s,\rho_s,\mathrm{AUC}_s$ over all held-out $\hat y_i$; residual quantiles $q_{0.1},q_{0.9}$
\EndFor
\State \Return means over seeds
\end{algorithmic}
\end{algorithm}

\paragraph{Controls and uncertainty.} On the same folds we score two baselines: the median training label, and the median
label of the training genomes that share the closest rank with the test genome (genus, family, order, class, phylum or domain),
which is what a reader could do with the taxonomy and the training table alone. Intervals are 95\% bootstrap intervals over
genera, or over the held-out families or orders (\NBootDraws draws), since species of one genus are not independent, with each statistic averaged over seeds in every draw. As a negative control the model is
trained on labels shuffled across genomes; it should then be no better than the median. Three robustness checks use the deployed settings over \TenNSeeds splits: the model against the relatives' median, the model with
gRodon's estimate (whose coefficients were fitted on data that include our species) replaced by a regression of the log
doubling time on gRodon's codon statistics fitted on the training rows of each split, and the model without the labels of
studies the audit flagged but kept. Each input group was removed in
turn, and rejected inputs added in turn, on the same genomes and folds over \NSeedsAblation seeds with the deployed settings
and each genome entered once (paired Wilcoxon tests across seeds).

\paragraph{Metagenome-assembled genomes.} For GEM genomes whose GTDB species has a measured growth rate, that species was
held out of training, and the prediction from each assembled genome was compared with the prediction from the isolate
genome and with the measured value (\MagNGenomes genomes, \MagNSpecies species, \NSeedsMag seeds).

\paragraph{Released tools.} gRodon \citep{weissman2021} and Phydon \citep{xu2025} were run as released on the same genes
and species-held-out folds. Phydon's distance to the nearest training relative still counted the held-out species itself,
so before each fold we removed that fold's species from Phydon's training table, without changing its code. Both tools take
the measured growth temperature; our model is compared without it and with it added as an input, built inside each split
as in Algorithm~\ref{alg:nested}. Per-genome absolute log errors, averaged over seeds, are compared with a paired Wilcoxon
test. The comparison was repeated with orders held out, removing every training species of the held-out orders from Phydon's
table (GTDB orders, matched by accession or genus for \PhydonOrderMapped\% of its species).

\paragraph{External test.} We searched PubMed abstracts and open full texts of species descriptions for doubling times
stated as the fastest, optimal or maximal growth in pure culture, and kept a value only if the species' NCBI taxonomy
identifier is absent from the compilation of \citet{madin2020}, the species is not in our labels, a genome assembly exists, and
its GTDB species has no training genome. Inputs were computed with the same code as for the labeled genomes, and the model was
trained once on all labeled genomes (\NSeedsExternal seeds averaged). As for GEM genomes, it predicts the doubling time at
its own predicted growth temperature, which is compared with the value as stated.

\subsection{Testing the low-energy explanation of slow growth}
Each lifestyle and energy input was tested against $y$ with growth temperature and domain as covariates, one row per
species group, by phylogenetic generalized least squares \citep{pagel1999,freckleton2002},
\begin{equation}
\mathbf{y}=\mathbf{Z}\boldsymbol\beta+\boldsymbol\varepsilon,\qquad
\boldsymbol\varepsilon\sim\mathcal{N}\big(\mathbf{0},\sigma^2[\lambda\mathbf{C}+(1-\lambda)\operatorname{diag}(\mathbf{C})]\big),
\end{equation}
where $\mathbf{C}$ holds the shared branch lengths on the GTDB trees and Pagel's $\lambda$ is estimated by maximum
likelihood. We also compared carriers and non-carriers within families and modeled the chance of $d>24$\,h by logistic
regression with \KEigen phylogenetic eigenvectors as covariates \citep{dinizfilho1998}. P values were adjusted for the
number of inputs \citep{benjamini1995}. The results were checked by leaving out one lineage or one species at a time,
excluding the literature strains, restricting to species at their own GTDB representative, and fixing $\lambda$ at 1.

\subsection{Application to the GEM catalog}
We predicted all \NGEMHQ high-quality genomes \citep{bowers2017} of the Genomes from Earth's Microbiomes catalog
\citep{nayfach2021} (\NGEMTotal genomes in total). Their GTDB taxonomy comes from the catalog and uses an older release, so
a renamed taxon looks more distant than it is and the novelty level is a worst case. Before predicting, the run stops if
the mean of any phylogeny or 16S input over the catalog genomes differs from the training mean by more than \DriftSD
training standard deviations, which guards against inputs computed in a different way.
